


\section{Setting Goals}

\begin{frame}{Data Structure of Tasks}
    
    \mytext{
    
     IMPORTANT: 
     
     \begin{itemize}
        \item a (result $\equiv$ milestone) can be checked!
        \item a (result $\equiv$ milestone) is an interface!
        \item a (result $\equiv$ milestone) leads to another task!
        \item all required resources are themselves dependencies and could in theory be modelled as tasks as well!
        \item however, cost, man hours, duration etc. are KPIs
        \item[] (Key Performance Indicators).
     \end{itemize}
    }
\end{frame}






\begin{frame}{Every Task Results in a Milestone}
    \mycomment{
    \textbf{Never} define a task without a clearly defined result, which can be checked, and which triggers the handover and closure of that task.
    }
    \showmytikz{../shared/tikz/tikz-PM-taskmilestone}
\end{frame}



% \begin{frame}{How to create a WBS}
    
%     \mytext{
   
%         First, with expert knowledge...
   
%         \begin{enumerate}
%             \item brainstorm what needs to be done!
%             \item group these tasks
%             \item[$\rightarrow$] flow chart, tree or mindmap format
%         \end{enumerate}
        
%     }
% \end{frame}

% \begin{frame}{How to create a WBS}
    
%     \mytext{
        
%         Then,for each task...
        
%         \begin{enumerate}
%             \item define goal ($\equiv$ result, milestone)
%             \item determine duration
%             \item determine requirements
%             \begin{itemize}
%                 \item dependency on other milestones
%                 \item resources required
%             \end{itemize}
%             \item assign responsibility
%             \begin{itemize}
%                 \item identify required competency
%             \end{itemize}
%         \end{enumerate}
%     }
% \end{frame}

% \begin{frame}{Flow Charts, Lists, Tree Chart, Gannt Diagramm}
    
%     \showmytikz{../shared/tikz/tikz-wbs-tree-gantt-list}
        
% \end{frame}


    % Considering your PM tool,
    % \begin{itemize}
    %     \item How easy is it to create a WBS?
    %     \item How easy is it to create a list of tasks?
    %     \item What data structure / monitoring does it permit?
    %     \item Is it easy to swap between representations?
    %     \item Is it easy to work with for a group?
    %     \item Does it integrate with other stuff?
    % \end{itemize}
    
    % \framebreak
 
    % Considering your project
    %     \begin{enumerate}
    %         \item What main work packages can you make out?
    %         \item What main milestones can you make out?
    %         \item What responsible persons do you see?
    %     \end{enumerate}

    
%     }
% \end{frame}
    


% \begin{frame}{Phase Oriented Project Management}
%     \resizebox{!}{0.7\textheight}{
%     \input{tikz/tikz-phase-oriented-project-management}
%     }
    
%     Source: \cite{schiersmannthiel}
% \end{frame}

% \begin{frame}{Typical Main Project Phases (Suchanek)}
%     \resizebox{!}{0.8\textheight}{
%         \input{tikz/tikz-chain-projectphases-details}
%     }

%     \framebreak
    
%     \resizebox{\textwidth}{!}{
%     \input{gantt-mains}
%     }
% \end{frame}



% \begin{frame}{Project Launch per Kerzner}

% \resizebox{!}{0.4\textheight}{
% \parbox{\textwidth}{
%     In the launch phase, customer (= project sponsor), project manager (pm), subject matter experts (sme's) and functional leads (leads) are involved. 
%     \begin{description}
%     \item[Phase 1] - Initialisation
%     \begin{itemize}
%         \item Kickoff-Meeting with customer
%     \end{itemize}
%     \item[Phase 2] - Preliminary planning
%     \begin{itemize}
%         \item PM, SMEs, Leads: Scope Management plan
%         \item Perliminary Scope Statement
%     \end{itemize}
%     \item[Phase 3] - Detailed planning
%     \begin{itemize}
%         \item PM, SMEs, Leads: \textbf{Work Breakdown Structure (WBS)} development
%         \item WBS dictionary
%     \end{itemize}
%     \item[Phase 4] Scope baseline
%     \begin{itemize}
%         \item Finalisation of governing documents\\ ("Lastenheft", product requirement documentation)
%     \end{itemize}
%     \end{description}
    
%     Source: \cite[p. 158]{KerznerHarold2017Pm:a}
% }
% }    
% \end{frame}

% \begin{frame}{The Objectives of a Project}

%     Well defined objectives are prerequisite. BUT
%     \begin{itemize}
%         \item objectives may be subject to change.
%         \item different stake holders may have different views.
%     \end{itemize}
    
    
%     Defined roles and institutions are much less easy to change.
    
% \end{frame}



% \begin{frame}{Project Flow Chart}

    
%     \input{tikz/tikz-projectflowchartone.tex}
    
%     \framebreak


%     \input{tikz/tikz-projectflowcharttwo.tex}

    
%     \framebreak
    
%     \input{tikz/tikz-projectflowchartthree.tex}
    

    
    
% \end{frame}







% \begin{frame}{Dimensions}

%     A project
%     \begin{itemize}
%         \item induces change
%         \item involves people 
%         \item requires resources
%         \item has external effects
%     \end{itemize}

%     A project has
%     \begin{itemize}
%         \item technical dimension (`technisches Angebot')
%         \item business dimension ('kaufmännisches Angebot')
%         \item human dimension (?)
%     \end{itemize}

%     \framebreak

%     Every project has multiple dimensions:
%     \begin{itemize}
%             \item technical
%             \item legal
%             \item economic
%             \item social
%             \item ecological
%             \item IT, data
%             \item ...
%     \end{itemize}
    
%     \framebreak
    
%     The environment could be
%     \begin{itemize}
%         \item hostile or friendly
%         \item stable, or dynamic and reactive
%     \end{itemize}
%     % This can be true for
%     % \begin{itemize}
%     %     \item technology
%     %     \item political / legal situation
%     %     \item economy
%     %     \item market
%     % \end{itemize}
    
%     $\rightarrow$ Different management styles may be required
% \end{frame}





% \input{break}

%\subtitle{Projects, Tasks, Goals and Objectives}


%\frame{\titlepage}

% \begin{frame}{Recap}

%     \begin{itemize}%[<+->]
%         \item A project is unique
%         \item A project delivers change
%         \item A project ventures into the unknown
%         \item A project is supposed to deliver value
%         \item A project consists of tasks
%         \item Complex projects have sub-projects
%     \end{itemize}
    
%     % \framebreak
    
%     % \begin{enumerate}%[<+->]
%     %     \item Who wants the project to happen?
%     %     \item How will success be measured?
%     %     \item Where is the project located?
%     %     \item What is required?
%     %     \item How do we get it done?
%     %     \begin{itemize}
%     %         \item Personel
%     %         \item Material
%     %         \item Other stakeholders
%     %         \item $\dots$
%     %     \end{itemize}
%     % \end{enumerate}
% \end{frame}




% \begin{frame}{WBS Work Breakdown Structure}  % depends on: forest
%     \begin{forest}
%     for tree={
%         font=\footnotesize,
%         grow = east,
%         parent anchor = east,
%         child anchor = west,
%         forked edge,
%         inner sep=2pt
%             }
%     [Travel to Jerusalem
%         [arrive]
%         [start walking [face problems]]
%         [wake up]
%     ]
%     \end{forest}

    
% \end{frame}


% \begin{frame}{WBS Work Breakdown Structure}  % depends on: forest
%     \begin{forest}
%     for tree={
%         font=\footnotesize,
%         grow = east,
%         parent anchor = east,
%         child anchor = west,
%         forked edge,
%         inner sep=2pt
%             }
%     [Travel to Jerusalem
%         [Reviewing [create picture book][review]]
%         [Execution [go to airport][pack stuff]]
%         [Preparation [...][book hotel][book flight][check finances / start saving]]
%     ]
%     \end{forest}
    
        
%     \framebreak
    
%     \smartdiagram[sequence diagram]{preparation, execution, reviewing}
    
%     `We cross the bridge once we come to it' - philosophy
    
%     Foresight?
    
% \end{frame}


% \begin{frame}{WBS Hold a Conference}  % depends on: forest
%     \begin{forest}
%     for tree={
%         font=\footnotesize,
%         grow = east,
%         parent anchor = east,
%         child anchor = west,
%         forked edge,
%         inner sep=0.5pt
%             }
%     [Hold a conference
%         [Review]
%         [Execution [venue management][stage service][moderation][welcome participants] ]
%         [Physical planning [contracts][bookings] ]
%         [Content planning [curate: choose presenters][Invite][Book space][Book services]]
%         [Preparation [set up team][finance][plan][advertise][conceptualise]]
%     ]
%     \end{forest}
    
% \end{frame}



% \begin{frame}{WBS Expedition}  % depends on: forest
%     \begin{forest}
%     for tree={
%         font=\footnotesize,
%         grow = east,
%         parent anchor = east,
%         child anchor = west,
%         forked edge,
%         inner sep=0.5pt
%             }
%     [Expedition to Antarctica
%         [Review 
%             [publication]
%             [stakeholder report]
%             [closing]
%             ]
%         [Execution  
%             [scientific research]
%             [watch out weather etc.]
%             [navigate]
%             [set sail / logistics]
%         ]
%         [Preparation 
%             [schedule]
%             [formulate goal]
%             [plan expenses]
%             [find business case] 
%             [identify sponsors]
%             [find vessel]
%             [build up team]
%         ]
%     ]
%     \end{forest}
% \end{frame}




% \begin{frame}{Product Breakdown}
%     \begin{forest}
%     for tree={
%         font=\footnotesize,
%         grow = west,
%         parent anchor = west,
%         child anchor = east,
%         forked edge,
%         inner sep=0.5pt
%             }
%     [New Factory
%         [Managment documents
%             [Business case]
%             [Project governance report]
%             [Project closure report]
%         ]
%         [Legal documents
%             [Approved insurance documents]
%             [Usage certificate]
%         ]
%         [Factory requirements
%             [Use requirements document]
%             [Solution design document]
%         ]
%         [Operational factory
%             [Selected factory]
%             [Fitted internal walls]
%             [Installed furniture and fixtures]
%             [Installed IT and office equipment]
%             [Fitted factory machinery]
%             [Marketed factory]
%         ]
%         [Factory Employees]
%         [Launched Factory]
%         ]
%     \end{forest}
    
%     Source: \cite[p. 168]{Roberts}

% \end{frame}

\begin{frame}{Tasks, Work Packages, Milestones}
    \mycomment{
    So you must define results for your work packages and your tasks. You can then let the internal planning of each task (and work package) to whoever is responsible. 
    }
    
    \mytext{
    \begin{itemize}
        \item Projects are defined by outcomes
        \item Tasks / Work Packages too
        \item Outcomes are \textbf{milestones}
        \item Milestones are \textbf{interface}
    \end{itemize}
    }
\end{frame}

\begin{frame}{Goal Setting / Defining Milestones}
    
    \mytext{
    \begin{itemize}
        \item How to set clear goals
        \item Useful / measurable mile stones
        \item Measure quality, quantity, \textbf{time}
    \end{itemize}
    }
\end{frame}


% \begin{frame}{Product Flow Chart (part)}
%     \tikzstyle{decision} = [diamond, draw, fill=blue!20, 
%         text width=4.5em, text badly centered, node distance=2cm, inner sep=0pt]
%     \tikzstyle{block} = [rectangle, draw, fill=blue!20, 
%         text width=5em, text centered, rounded corners, minimum height=4em]
%     \tikzstyle{line} = [draw, -latex']
%     \tikzstyle{cloud} = [draw, ellipse,fill=red!20, node distance=3cm,
%         minimum height=2em]
        
%     \begin{tikzpicture}[scale=0.2, node distance = 2.2cm, auto]
%     \footnotesize
%         % Place nodes
%         \node [block] (case) {business case};
%         \node [block, below of=case] (gov) {governance document};
%         \node [block, below left of=gov] (use) {user requirements document};
%         \node [block, below right of=gov] (sol) {solution design document};
%         \node [block, below right of=use] (sel) {selected factory};
%         % Draw edges
%         \path [line] (case) -- (gov);
%         \path [line] (gov) -- (use);
%         \path [line] (gov) -- (sol);
%         \path [line] (use) -- (sel);
%         \path [line] (sol) -- (sel);
%     \end{tikzpicture}
    
%     Source: \cite[p. 169]{Roberts}
    
% \end{frame}

\begin{frame}{The Art of Setting Clear Goals}
    
    \mytext{
    \begin{itemize}
        \item Define precisely what you want
        \item ...and from whom
        \item ...and when
        \item ...and how you will check
        \item ...and what happens if...
    \end{itemize}
    }
\end{frame}

\begin{frame}{The Three Constraints}
    \showmytikz[caption={The Three Constraints (Planning)},label=fig:threeconstraints-planning]{../shared/tikz/tikz-threeconstraints}
    \mycomment{
    You can have it \textbf{cheap,} \textbf{fast, }or \textbf{good.}
    
    Choose two!
    
    \vfill
    \begin{itemize}
        \item Which is more `flexible': quality, time, costs?
        \item Is any one factor critical / not disputable?
        \item Do all stake holder see this in the same way?
        \item How do they interact?
    \end{itemize}
    }
\end{frame}

\begin{frame}{SMART Goals}
    \mycomment{
    Not only do you need to set goals, but these must also be clearly defined in a number of dimensions, and also be achievable altogether. 
    }
    \mytext{
    Goals must be \colorbox{yellow!40}{SMART}:
    
    \begin{description}
        \item[S] specific (simple, sensible, significant)
        \item[M] measurable (meaningful, motivating)
        \item[A] achievable (agreed, attainable)
        \item[R] relevant (reasonable, realistic and resourced, results-based)
        \item[T] time bound (time-based, time limited, time/cost limited, timely, time-sensitive)
    \end{description}
 
    
    % German `operational definiert'
    
    }
\end{frame}

\begin{frame}{Goal Setting for Success }
    \mycomment{
    Locke and Latham \cite{LOCKE1968157} have worked out that goals you set (for your employees) need to fulfil five critera for success. The goal must be stated clearly; it must contain a bit of a challenge; you must request commitment to this goal from the person you have set it to; you must take and provide feedback; and the task has to have the right amount of complexity: highly complex tasks require extra time and buffer. 
    }
    \mytext{
    Five goal-setting principles for success:
    
    \begin{enumerate}
        \item Clarity
        \item Challenge
        \item Commitment
        \item Feedback
        \item Task complexity
    \end{enumerate}
    
    \cite{LOCKE1968157}
        
    }
\end{frame}

\begin{frame}{Tasks, Milestones}
    
    \mytext{
    \begin{itemize}
        \item s.m.a.r.t. goals can be milestones
        \item tasks need planning also... or do they?
        \item task/milestone definition: carefully!
        \item is the milestone flexible in time, scope, costs?
        \item is there a way to review goals?
        \item whom is the responsibility attributed to?
    \end{itemize}
    }
\end{frame}




% \begin{frame}{Products vs.~Tasks and Time}
%     \begin{itemize}
%         \item Producs need to be produced
%         \item Production must be planned
%         \item Production takes time: tasks take time!
%         \item Move from a product view to a process view
%         \item Time is the most important \emph{influenceable }cost factor
%         \item How can we estimate time?
%     \end{itemize}
% \end{frame}

% \begin{frame}{Scope Creep}
%     \begin{itemize}
%         \item Sponsor realises he wants more
%         \item Management realises product needs more
%         \item Milking the sponsor or milking the mananger?
%         \item What is necessary?
%         \item How can this all be paid for?
%     \end{itemize}
% \end{frame}

% % \begin{frame}{Product Description}
% %     \begin{itemize}
% %         \item Title
% %         \item Purpose
% %         \item Composition
% %         \item Derivation
% %         \item Format
% %         \item Audience
% %         \item Quality criteria
% %         \item Quality method
% %     \end{itemize}
% % \end{frame}

% % \begin{frame}{The Three Constraints}
% %     \begin{tikzpicture}[scale=0.8]
% %         \coordinate (A) at (0,0);
% %         \coordinate (B) at (-3,-5);
% %         \coordinate (C) at (3,-5);
% %         \coordinate (X) at (0,-3.35);
% %         \draw (A) -- (B) node[midway,above,rotate=60] {Costs};
% %         \draw (B) -- (C) node[midway,below] {Quality/Performance/Technology};
% %         \draw (C) -- (A) node[midway,above,rotate=-60] {Time};
% %         \path[draw=blue, -stealth, rotate=90,
% %       postaction={decorate,decoration={text effects along path,
% %         text={\ Within good customer relations \ }, text align=center, 
% %         text effects/.cd, 
% %           text along path, 
% %           every character/.style={fill=white, yshift=-0.5ex}}}]
% %       (X) circle (4);
% %     \end{tikzpicture}
    
% %     Probably the most important slide.
% % \end{frame}

% \begin{frame}{Scope Freeze Milestones}
%     \begin{itemize}
%         \item A date where the scope cannot be altered
%         \item Typically, specifications are set
%         \item Also design freeze, customer approval freeze...
%         \item De-couple sponsor from product
%         \item De-couple project management from product?
%         \item There is agile!!
%     \end{itemize}
% \end{frame}

% \begin{frame}{MoScCoW - Scheme}
%     \begin{description}
%         \item[M] must have: non-negotiable
%         \item[S] should have: not vital, but adds significant value
%         \item[C] could have: nice to have but small impact when left out
%         \item[W] will not have: stuff without priority in the given time frame.
%     \end{description}    
% \end{frame}

\begin{frame}{Plan planning!}
    
    \mytext{
    Planning 
    \begin{itemize}
        \item creates a model or scenarios
        \item gathers information about possible futures
        \item is a process that takes time in itself
        \item may be self-referential and prone to loops
        \item aims at changing the future, creating value
        \item is necessarily simpler than reality
        \item will get more specific the closer the time
        \begin{itemize}
            \item have a very clear plan for tomorrow
            \item another for the near future
            \item another for the far future
            \item and values for orientation
        \end{itemize}
    \end{itemize}
    }
\end{frame}

% \begin{frame}{Famous Words}
% \begin{quote}
%     Culture eats strategy for breakfast.
    

% \end{quote}
%     \medskip
%     \flushright ---Attr.: Peter \noun{Drucker}
    
% \end{frame}



% \begingroup
%     \ifdefined\setbeamercolor\setbeamercolor{background canvas}{bg=violet!20}\fi
%     \begin{frame}{TASK}
    
%     \mytext{
%     Create a Gantt-Diagram for the Project Loft Bed (slide~\ref{sl:loftbed}), defining tasks, milestones etc.
%         \begin{enumerate}
%             \item How long will the project take?
%             \item What are the most important tasks, and why?
%             \item Which tasks can be done in parallel?
%             \item How many work hours will you need?
%             \item Which skill sets will you need?
%         \end{enumerate}
    
%     }
% \end{frame}
% \endgroup
 
 
%  \begin{frame}{Your example projects!}
%     In groups:
%     \begin{itemize}
%         \item Present each other one of your example projects
%         \item Who wanted what?
%         \item Was it delivered in time, budget, quality?
%     \end{itemize}
    
%     Share one example!
%     \parbox{8cm}{    
%     \centering
%     \clock[0.9]{11}{45}{5}
%     }
% \end{frame}
